\chapter{Discussion}

\section{Interpretation of the primary Results}

Construction strategy clearly influenced the number of broad-phase bounding-volume checks.
Bottom-up produced the lowest mean in 17 out of 18 parameter combinations, whereas top-down produced the highest in all 18.
The analysis of all scenes, shows this pattern is visible in 885 out of 900 scenes, so the ordering did not result from few extreme results.

The construction criteria provide one plausible explanation.
Bottom-up selects the pair for the active-node whose merged AABB has the smallest area, leading to local clusters forming.
Top-down splits the longera axis at the median, leading to a more balanced tree without taking overlap or area into account.
Incremental insetion considers AABB-area growth but searches only one specific path.

As sibling overlap and AABB area by depth were not recorded, this interpretation remains speculative. 
Therefore no single structural property can be identified as the cause.


\section{Scope and Limitations}

The conclusion applies to the three concreate implemented algorithms and heuristics used in this experiment.
It does not establish that bottom-up construction is generally superior to every top-down or incremental method.
A top-down method using e.g. a surface-area heursitic could therefore produce a different results from the median split evaluated here.

The chosen metrics are also proxies for traversal cost only rather than measurements of execution time. 
A Bounding volume construction method may construct a hierarchy that is fast to traverse, but takes significatnly longer to construct, thereby reducing the advantage gained.
This trade-off is acceptable when a hierarchy is traversed many times, but total performance cannot be assessed without construction time and the expected number of queires to the hierarchy taken into account.

Further limitations come from the scene model.
The experiment uses uniformly generated static two-dimensional scenes, equal-sized circles, AABBs and the tested ranges of 25 to 1000 particles and 5\% to 20\% coverage.
The results cannot directly be transfered to different scenarios.
The 50 datasets per parameter combination show that the observed ordering is robust for the investigated model.


\newpage
